\chapter{Algorytmiczne przeszukiwanie grafów: BFS i DFS}
\label{ch:tg-bfs-dfs}

Rozdziały \ref{ch:tg-podstawy}--\ref{ch:tg-skojarzenia}
wprowadziły ścieżki, odległość, drzewa rozpinające, cykle i spójność.
W rozdziale \ref{ch:tg-drzewa} wspomnieliśmy o przeszukiwaniu wszerz
i w głąb, a w części poświęconej orientacjom
wykorzystaliśmy własność drzewa DFS.
Teraz podamy oba algorytmy, niezmienniki zapewniające ich poprawność
oraz dokładne oszacowania kosztu.

\section{Model obliczeń i reprezentacja grafu}
\label{sec:tg-algo-model}

W całym rozdziale $G=(V,E)$ jest \emph{skończonym grafem prostym
nieskierowanym}, $n=|V|$, $m=|E|$, chyba że jawnie mówimy
o grafie skierowanym $D=(V,A)$.
Dla $v\in V$ piszemy $N(v)=\{w\in V:\{v,w\}\in E\}$.
Zakładamy, że wierzchołki mają identyfikatory, dzięki którym
odczyt i zapis wartości tablicy indeksowanej przez $V$
kosztuje $O(1)$.
Graf przechowujemy jako \emph{listy sąsiedztwa}:
dla każdego $v$ dostępna jest lista wszystkich elementów $N(v)$.
Utworzenie takiej reprezentacji z listy krawędzi trwa $O(n+m)$
i zajmuje $O(n+m)$ pamięci.
Każda krawędź grafu nieskierowanego pojawia się na dwóch listach.
Kolejność sąsiadów może wpływać na otrzymane drzewo przeszukiwania,
ale nie na prawdziwość udowodnionych własności.

\begin{definicja}[Odkrycie, przetworzenie i rodzic]
\label{def:tg-algo-discovery}
Wierzchołek jest \emph{nieodkryty}, zanim algorytm po raz pierwszy
do niego dotrze, oraz \emph{odkryty} od tej chwili.
Tablica rodziców $p$ przypisuje każdemu wierzchołkowi odkrytemu
z innego wierzchołka $v$ jego poprzednika $p(v)$.
Dla korzenia przeszukiwania oraz wierzchołków nieosiągalnych
przyjmujemy $p(v)=\bot$.
Krawędzie $\{p(v),v\}$, dla których $p(v)\neq\bot$,
nazywamy \emph{krawędziami drzewa przeszukiwania}.
\end{definicja}

W pseudokodzie \texttt{Q.push} dopisuje element na koniec kolejki,
\texttt{Q.pop} usuwa element z początku,
a \texttt{for w in N(v)} przegląda sąsiadów w ustalonej kolejności.
Dla obliczeń złożoności przyjmujemy jednostkowy koszt tych operacji
oraz dodawania i porównywania liczb całkowitych z zakresu $0,\ldots,n$.
Nie używamy reprezentacji macierzą sąsiedztwa przy wyprowadzaniu
granic $O(n+m)$.

\section{Przeszukiwanie wszerz (BFS)}
\label{sec:tg-algo-bfs}

Niech $s\in V$. Przeszukiwanie wszerz (BFS, od ang.
\emph{breadth-first search}) odwiedza wierzchołki kolejnymi
warstwami odległości od $s$.
Tablica $d$ przechowuje długość drogi odkrytej przez algorytm.

\begin{definicja}[Algorytm BFS]
\label{def:tg-algo-bfs}
Początkowo dla każdego $v\in V$ ustawiamy
$\mathrm{seen}[v]=\mathrm{false}$, $d[v]=\infty$,
$p[v]=\bot$. Następnie wykonujemy kroki:
\begin{enumerate}[label=(\arabic*)]
\item Ustawiamy $\mathrm{seen}[s]=\mathrm{true}$, $d[s]=0$
      i umieszczamy $s$ w pustej kolejce $Q$.
\item Dopóki $Q$ nie jest pusta, usuwamy z jej początku
      wierzchołek $v$.
\item Dla kolejnych $w\in N(v)$, dla których
      $\mathrm{seen}[w]=\mathrm{false}$, ustawiamy
      $\mathrm{seen}[w]=\mathrm{true}$, $p[w]=v$,
      $d[w]=d[v]+1$, po czym dopisujemy $w$ na koniec $Q$.
\item Kończymy, gdy kolejka jest pusta.
\end{enumerate}
Ważne: wierzchołek oznaczamy jako odkryty już przy
\emph{dopisaniu} do kolejki, nie dopiero podczas usuwania.
\end{definicja}

\begin{lemat}[Niezmiennik kolejki]
\label{lem:tg-bfs-queue}
W dowolnej chwili algorytmu BFS, gdy kolejka nie jest pusta,
wartości $d$ kolejnych jej elementów są niemalejące,
a różnica wartości dla ostatniego i pierwszego elementu
jest co najwyżej jeden. Ponadto każdy wierzchołek zostaje
dopisany do kolejki co najwyżej raz.
\end{lemat}
\begin{proof}
Na początku kolejka składa się wyłącznie z $s$, więc
pierwsza własność jest spełniona. Załóżmy ją przed
usunięciem elementu $v$ o wartości $k=d[v]$.
W pozostałej części kolejki wartości są równe $k$
lub $k+1$, przy czym wszystkie $k$ poprzedzają $k+1$.
Każdy nowo odkryty sąsiad $v$ otrzymuje wartość $k+1$
i trafia na koniec kolejki. Powstała kolejka
ma nadal wartości niemalejące i najwyżej dwie
kolejne wartości. Indukcja po liczbie usunięć
daje tezę, także gdy kolejka przejściowo staje się pusta.
Oznaczanie wierzchołka przy pierwszym dopisaniu sprawia,
że żadne późniejsze przeglądanie krawędzi nie dopisze go
powtórnie.
\end{proof}

\begin{twierdzenie}[Poprawność i optymalność BFS]
\label{thm:tg-bfs-correct}
Niech $G=(V,E)$ będzie skończonym grafem prostym
nieskierowanym, $s\in V$, a $d,p$ wynikiem BFS z $s$.
Wówczas dla każdego $v\in V$ zachodzą następujące własności:
\begin{enumerate}[label=(\roman*)]
\item $v$ zostaje odkryty wtedy i tylko wtedy, gdy istnieje
      droga z $s$ do $v$;
\item $d[v]=d_G(s,v)$, z przyjęciem wartości $\infty$
      dla wierzchołków nieosiągalnych;
\item gdy $v\neq s$ jest osiągalny, kolejne zastosowania $p$
      odtwarzają drogę z $s$ do $v$ długości $d_G(s,v)$;
\item krawędzie rodziców osiągalnych wierzchołków tworzą
      drzewo rozpinające składowej spójnej zawierającej $s$.
\end{enumerate}
\end{twierdzenie}
\begin{proof}
Każdy odkryty wierzchołek inny niż $s$ ma rodzica
odkrytego wcześniej, a jego wartość $d$ jest o jeden
większa od wartości rodzica. Podążając za rodzicami,
otrzymujemy zatem po skończonej liczbie kroków $s$.
Odwrócenie tego ciągu daje drogę długości $d[v]$.
W szczególności $d_G(s,v)\leq d[v]$.

Pokażemy przez indukcję po $k\geq0$, że każdy wierzchołek
o prawdziwej odległości $k$ zostaje odkryty z wartością
$d[v]\leq k$. Dla $k=0$ jest to $s$.
Niech $d_G(s,v)=k+1$ i niech $u$ będzie przedostatnim
wierzchołkiem na najkrótszej drodze z $s$ do $v$.
Z założenia indukcyjnego $u$ zostaje odkryty i
$d[u]\leq k$. Algorytm usunie w końcu $u$ z kolejki,
ponieważ każdy odkryty wierzchołek jest usuwany dokładnie
raz, a kolejka nie może nieskończenie rosnąć.
Jeśli $v$ nie zostało jeszcze odkryte, przeglądanie
$\{u,v\}$ ustali $d[v]=d[u]+1\leq k+1$.
Jeśli $v$ odkryto wcześniej z innym rodzicem,
niezmiennik kolejki z lematu
\ref{lem:tg-bfs-queue} zapewnia, że wcześniejsze
odkrycie nie nastąpiło z wartością większą niż
$d[u]+1$: przed przetworzeniem $u$ algorytm nie może
przetworzyć elementu o wartości większej niż $d[u]$,
a każdy nowy element otrzymuje wartość o jeden większą.
Zatem zawsze $d[v]\leq k+1$.

Otrzymujemy $d[v]=d_G(s,v)$ dla wszystkich wierzchołków
osiągalnych. Każdy wierzchołek na drodze z $s$
zostaje odkryty na mocy powyższej indukcji;
w drugą stronę konstrukcja rodziców gwarantuje
istnienie drogi do odkrytego wierzchołka.
Stąd także (i), (ii) i (iii).
Ponieważ wartości $d$ ściśle maleją wzdłuż krawędzi
do rodzica, zbiór tych krawędzi nie zawiera cyklu,
a każdy osiągalny wierzchołek łączy się z $s$.
Jest więc drzewem na całej składowej $s$.
\end{proof}

\begin{przyklad}[Warstwy BFS i drzewo najkrótszych dróg]
\label{ex:tg-bfs-layers}
Niech $V=\{s,a,b,c,d,e,f\}$ oraz
$E=\{sa,sb,ac,ad,bd,be,cd,df,ef\}$, gdzie
$xy$ oznacza krawędź $\{x,y\}$.
Przy kolejności sąsiadów alfabetycznej, z $s$ jako
pierwszym elementem kolejki, otrzymujemy
$d[s]=0$, $d[a]=d[b]=1$,
$d[c]=d[d]=d[e]=2$ i $d[f]=3$.
Możliwe drzewo BFS ma krawędzie
$sa,sb,ac,ad,be,df$.
Krawędzie $bd,cd,ef$ nie należą do tego drzewa;
nie wpływa to na poprawność obliczonych odległości.
\end{przyklad}

\begin{figure}[htbp]
\centering
\begin{tikzpicture}[scale=.9]
\node[tgBlueVertex] (s) at (0,0) {$s$};
\node[tgVertex] (a) at (2,1.35) {$a$};
\node[tgVertex] (b) at (2,-1.35) {$b$};
\node[tgVertex] (c) at (4,1.55) {$c$};
\node[tgVertex] (d) at (4,0) {$d$};
\node[tgVertex] (e) at (4,-1.55) {$e$};
\node[tgVertex] (f) at (6,0) {$f$};
\draw[tgTree] (s)--(a) (s)--(b) (a)--(c) (a)--(d)
 (b)--(e) (d)--(f);
\draw[tgGhost] (b)--(d) (c)--(d) (e)--(f);
\node[font=\footnotesize,text=tgNavy] at (0,-.6) {$d=0$};
\node[font=\footnotesize,text=tgNavy] at (2,-2.2) {$d=1$};
\node[font=\footnotesize,text=tgNavy] at (4,-2.2) {$d=2$};
\node[font=\footnotesize,text=tgNavy] at (6,-.6) {$d=3$};
\end{tikzpicture}
\caption{Przykładowe warstwy BFS od $s$: krawędzie drzewa
oznaczono linią ciągłą, pozostałe -- przerywaną.}
\label{fig:tg-bfs-layers}
\end{figure}

\begin{stwierdzenie}[Złożoność BFS]
\label{prop:tg-bfs-complexity}
Dla grafu prostego zadanego listami sąsiedztwa BFS
wykonuje $O(n+m)$ operacji i zużywa $O(n)$
pamięci pomocniczej poza wejściową reprezentacją grafu.
\end{stwierdzenie}
\begin{proof}
Inicjalizacja tablic kosztuje $O(n)$.
Każdy wierzchołek zostaje dopisany i usunięty
co najwyżej raz. Lista sąsiedztwa $v$
jest przeglądana tylko podczas usuwania $v$,
zatem łączna liczba odczytów wpisów nie przekracza
$\sum_{v\in V}d_G(v)=2m$, także dla grafu niespójnego.
Wszystkie pozostałe czynności przy jednym wpisie
zajmują $O(1)$. Tablice i kolejka przechowują po
co najwyżej $n$ elementów.
\end{proof}

\section{Zastosowania BFS: składowe i dwudzielność}
\label{sec:tg-bfs-applications}

\begin{twierdzenie}[Składowe spójne w czasie liniowym]
\label{thm:tg-components-algo}
Dla dowolnego skończonego grafu prostego $G$ można
wyznaczyć wszystkie składowe spójne w czasie $O(n+m)$,
uruchamiając BFS kolejno z każdego jeszcze
nieodkrytego wierzchołka i nie zerując tablicy
\texttt{seen} pomiędzy uruchomieniami.
\end{twierdzenie}
\begin{proof}
Z twierdzenia \ref{thm:tg-bfs-correct} jedno uruchomienie
odkrywa dokładnie jedną składową. Kolejny korzeń
należy do innej składowej, a po zakończeniu
nie zostaje żaden nieodkryty wierzchołek.
Listy sąsiedztwa każdej składowej przeglądamy
dokładnie raz, więc sumaryczny koszt wynosi
$O(n+\sum_v d_G(v))=O(n+m)$.
\end{proof}

\begin{twierdzenie}[Test dwudzielności]
\label{thm:tg-bfs-bipartite}
Dla skończonego grafu prostego $G$ uruchommy BFS
w każdej składowej i przypiszmy odkrytemu
wierzchołkowi kolor $d[v]\bmod2$.
Graf jest dwudzielny wtedy i tylko wtedy,
gdy dla każdej krawędzi $\{u,v\}$ kolory
końców są różne. Sprawdzenie trwa $O(n+m)$.
\end{twierdzenie}
\begin{proof}
Jeżeli test nie wykrywa krawędzi o końcach tego
samego koloru, podział na oba zbiory kolorów
spełnia definicję grafu dwudzielnego.
Jeżeli istnieje krawędź $uv$ o równych kolorach,
drogi w drzewie BFS od korzenia do $u$ i $v$
mają długości tej samej parzystości.
Niech $z$ będzie ich ostatnim wspólnym wierzchołkiem.
Połączenie rozłącznych odcinków dróg od $z$ do
$u$ oraz od $z$ do $v$ z krawędzią $uv$
tworzy cykl długości
$d[u]+d[v]-2d[z]+1$, a więc nieparzystej
(zauważmy, że $uv$ nie jest krawędzią drzewa).
Graf zawierający cykl nieparzysty nie jest dwudzielny,
bo na takim cyklu kolory musiałyby naprzemiennie się
zmieniać i nie mogłyby wrócić do wyjściowego.
Koszt badania krawędzi jest liniowy jak dla BFS.
\end{proof}

\section{Przeszukiwanie w głąb (DFS)}
\label{sec:tg-algo-dfs}

DFS (\emph{depth-first search}) bada jedną gałąź
aż do wyczerpania możliwości, po czym wraca
do najbliższego wierzchołka mającego
nieodkrytego sąsiada. Wersję rekurencyjną
można zastąpić wersją ze stosem, która
przechowuje także bieżącą pozycję
w liście sąsiedztwa każdego aktywnego wierzchołka.

\begin{definicja}[Algorytm DFS z czasami wejścia i wyjścia]
\label{def:tg-algo-dfs}
Wszystkie wierzchołki są początkowo białe,
$p[v]=\bot$, a licznik czasu $t=0$.
Wykonujemy procedurę \texttt{DFS(v)}:
\begin{enumerate}[label=(\arabic*)]
\item Kolorujemy $v$ na szaro,
      zwiększamy $t$ o jeden i zapisujemy
      $\mathrm{tin}[v]=t$.
\item Dla każdego sąsiada $w\in N(v)$:
      jeśli $w$ jest biały, ustawiamy $p[w]=v$
      i wywołujemy \texttt{DFS(w)}.
\item Po przeglądnięciu całej listy $N(v)$
      kolorujemy $v$ na czarno, zwiększamy
      $t$ o jeden i zapisujemy $\mathrm{tout}[v]=t$.
\end{enumerate}
W celu przeszukania całego grafu wywołujemy
\texttt{DFS(v)} z każdego jeszcze białego
wierzchołka. Szare wierzchołki odpowiadają
dokładnie aktywnemu stosowi rekurencji.
\end{definicja}

\begin{twierdzenie}[Las DFS i przedziały czasowe]
\label{thm:tg-dfs-intervals}
Dla skończonego grafu prostego nieskierowanego
wynik DFS spełnia:
\begin{enumerate}[label=(\roman*)]
\item krawędzie rodziców tworzą las rozpinający $G$,
      po jednym drzewie na każdą składową spójną;
\item dla różnych wierzchołków $u,v$ zachodzi
      dokładnie jedna z relacji:
      przedziały $[\mathrm{tin}[u],\mathrm{tout}[u]]$
      i $[\mathrm{tin}[v],\mathrm{tout}[v]]$ są rozłączne
      albo jeden jest właściwie zawarty w drugim;
\item $u$ jest przodkiem $v$ w lesie DFS
      (dopuszczamy $u=v$) wtedy i tylko wtedy,
      gdy przedział $v$ zawiera się
      w przedziale $u$;
\item każda krawędź niedrzewowa łączy wierzchołek
      z jego właściwym przodkiem w lesie DFS.
\end{enumerate}
\end{twierdzenie}
\begin{proof}
Każdy wierzchołek otrzymuje rodzica tylko w chwili
pierwszego odkrycia. Rodzic jest odkryty wcześniej
i pozostaje szary do końca wywołania potomka.
Nie może więc powstać cykl skierowany ku rodzicom.
Z dowolnego korzenia DFS docieramy do wszystkich
osiągalnych wierzchołków: gdyby pierwszy wierzchołek
nieodkryty na drodze wychodzącej od korzenia
istniał po zakończeniu procedury, jego poprzednik
musiałby wcześniej przejrzeć odpowiednią krawędź
i go odkryć. Stąd (i).

Wykonanie rekurencyjne ma własność zagnieżdżenia:
wywołanie potomka kończy się, zanim skończy się
wywołanie jego przodka. Wywołania w różnych
gałęziach nie nakładają się w czasie.
Wynikają stąd (ii) i (iii), gdyż każdy czas
wejścia i wyjścia jest odrębny.

Dla (iv) załóżmy, że końce krawędzi $uv$
nie są powiązane relacją przodka.
Z (ii) ich przedziały byłyby rozłączne;
niech bez straty ogólności $u$ kończy się
przed odkryciem $v$. Podczas wykonywania
\texttt{DFS(u)} przeglądamy jednak krawędź
$uv$. Jeśli $v$ był wówczas biały,
procedura odkryłaby go jako potomka $u$,
wbrew założeniu. Jeśli nie był biały,
jego wywołanie nie mogłoby rozpoczynać się
dopiero po zakończeniu $u$.
Sprzeczność dowodzi (iv).
\end{proof}

\begin{stwierdzenie}[Złożoność DFS]
\label{prop:tg-dfs-complexity}
DFS dla list sąsiedztwa działa w czasie $O(n+m)$
i zajmuje $O(n)$ pamięci pomocniczej,
licząc stos rekurencji.
\end{stwierdzenie}
\begin{proof}
Każdy wierzchołek rozpoczyna i kończy wywołanie
dokładnie raz, a każdą listę sąsiedztwa
przeglądamy najwyżej raz. Łącznie odczytujemy
$2m$ wpisów i wykonujemy $O(n)$ aktualizacji
kolorów, rodziców i czasów. Głębokość stosu
jest co najwyżej $n$, ponieważ w jednym
momencie aktywne wywołania odpowiadają
różnym wierzchołkom.
\end{proof}

\section{Mosty i wykrywanie cykli metodą DFS}
\label{sec:tg-algo-low}

Dla ustalonego drzewa DFS zakorzenionego w $r$
oznaczmy przez $T_v$ zbiór wierzchołków
poddrzewa o korzeniu $v$. Dla $v\in V$
zdefiniujmy $\mathrm{low}[v]$ jako minimum
$\mathrm{tin}[x]$ po wszystkich wierzchołkach
$x$, do których można dotrzeć z pewnego
$y\in T_v$ po co najwyżej jednej krawędzi
niedrzewowej biegnącej ku przodkowi,
dopuszczając również drogę długości zero
z $v$ do $v$ oraz krawędzie drzewa
prowadzące w dół do $y$.
W szczególności $\mathrm{low}[v]\leq\mathrm{tin}[v]$.
Równoważnie obliczamy to minimum rekurencyjnie
w czasie kończenia procedury DFS:
\[
\mathrm{low}[v]=\min\left(
\{\mathrm{tin}[v]\}\cup
\{\mathrm{tin}[w]:vw\in E,\ w\text{ jest właściwym przodkiem }v\}
\cup
\{\mathrm{low}[w]:p[w]=v\}
\right).
\]
Dla grafu \emph{prostego} przy przeglądaniu
krawędzi wykluczamy jedyną krawędź do rodzica.
Dla multigrafu trzeba odróżniać identyfikatory
krawędzi równoległych; pominięcie wszystkich
krawędzi do rodzica dałoby niepoprawny wynik.

\begin{twierdzenie}[Kryterium mostu metodą DFS]
\label{thm:tg-dfs-bridge}
Niech $G=(V,E)$ będzie skończonym grafem prostym,
a $p,\mathrm{tin},\mathrm{low}$ danymi z DFS.
Dla krawędzi drzewa $\{v,w\}$, dla której
$p[w]=v$, następujące warunki są równoważne:
\begin{enumerate}[label=(\roman*)]
\item $\{v,w\}$ jest mostem;
\item $\mathrm{low}[w]>\mathrm{tin}[v]$.
\end{enumerate}
Każdy most jest krawędzią drzewa DFS.
Wszystkie mosty można znaleźć w czasie $O(n+m)$.
\end{twierdzenie}
\begin{proof}
Jeżeli $\{v,w\}$ jest krawędzią niedrzewową,
łączy wierzchołek z przodkiem na mocy
\ref{thm:tg-dfs-intervals}(iv). Razem z drogą
drzewową tworzy cykl, więc nie jest mostem.

Rozważmy teraz krawędź drzewową $\{v,w\}$,
$p[w]=v$. Jedyną krawędzią drzewa łączącą
$T_w$ z $V\setminus T_w$ jest $\{v,w\}$.
Na mocy \ref{thm:tg-dfs-intervals}(iv)
każda inna krawędź przekraczająca ten przekrój
jest niedrzewowa i łączy potomka $w$
z $v$ albo z przodkiem $v$.
Taka krawędź istnieje dokładnie wtedy, gdy
$\mathrm{low}[w]\leq\mathrm{tin}[v]$.
Jeżeli istnieje, po usunięciu $\{v,w\}$
nadal mamy połączenie poddrzewa $T_w$
z resztą składowej; krawędź nie jest mostem.
Jeśli nie istnieje, usunięcie $\{v,w\}$
odcina całe $T_w$, więc jest mostem.
Aktualizacje \texttt{low} wymagają stałej
pracy podczas analizy każdej krawędzi
i każdego wierzchołka, dlatego zachowują
złożoność DFS.
\end{proof}

\begin{przyklad}[Mosty i liczby \texttt{low}]
\label{ex:tg-dfs-low}
Rozważmy graf o wierzchołkach $\{s,a,b,c,d\}$
i krawędziach $sa,ab,bc,ca,bd$.
Przy kolejności DFS $s,a,b,c,d$ mamy
$\mathrm{tin}[s]=1$, $\mathrm{tin}[a]=2$,
$\mathrm{tin}[b]=3$, $\mathrm{tin}[c]=4$
oraz $\mathrm{tin}[d]=6$.
Krawędź $ca$ daje
$\mathrm{low}[c]=2$, a stąd
$\mathrm{low}[b]=\mathrm{low}[a]=2$.
Wierzchołek $d$ jest liściem, więc
$\mathrm{low}[d]=6$.
W konsekwencji mostami są dokładnie
$sa$ i $bd$; krawędzie cyklu $abc$ mostami
nie są.
\end{przyklad}

\begin{stwierdzenie}[Wykrywanie cyklu]
\label{prop:tg-dfs-cycle}
Skończony graf prosty nieskierowany zawiera
cykl wtedy i tylko wtedy, gdy DFS znajdzie
krawędź do szarego wierzchołka różnego
od rodzica bieżącego wierzchołka.
Test można wykonać w czasie $O(n+m)$.
\end{stwierdzenie}
\begin{proof}
Taka krawędź łączy wierzchołek z jego
właściwym przodkiem i wraz z drogą
drzewową tworzy cykl. Jeśli graf ma cykl,
nie może być lasem; krawędzie drzewa DFS
tworzą las, zatem istnieje krawędź
niedrzewowa. Z twierdzenia
\ref{thm:tg-dfs-intervals}(iv) ma ona
końce w relacji przodka. Przeglądana
od potomka do przodka trafia na szary
wierzchołek różny od rodzica.
\end{proof}

\section{DFS w grafach skierowanych: porządek topologiczny}
\label{sec:tg-algo-dag}

Niech $D=(V,A)$ będzie skończonym grafem
skierowanym w znaczeniu rozdziału
\ref{ch:tg-euler}. W procedurze DFS przeglądamy
wyłącznie łuki wychodzące.
\emph{Skierowanym cyklem} nazywamy
ciąg różnych wierzchołków
$v_0,\ldots,v_{k-1}$ dla $k\geq2$
z łukami $v_i\to v_{i+1}$ oraz
$v_{k-1}\to v_0$; pętla, jeśli jest
dozwolona, również wyklucza acykliczność.
\emph{Porządek topologiczny} to liniowe
uszeregowanie wierzchołków, w którym
każdy łuk $u\to v$ biegnie od
wcześniejszego do późniejszego wierzchołka.

\begin{twierdzenie}[Rozpoznawanie DAG i sortowanie topologiczne]
\label{thm:tg-dfs-dag}
Niech $D=(V,A)$ będzie skończonym grafem
skierowanym bez pętli.
Wówczas $D$ jest acykliczny wtedy i tylko
wtedy, gdy podczas DFS nie wystąpi łuk
od bieżącego wierzchołka do szarego
wierzchołka. Jeśli $D$ jest acykliczny,
kolejność malejących wartości
$\mathrm{tout}$ jest porządkiem
topologicznym. Obie operacje mają
złożoność $O(|V|+|A|)$ dla list
sąsiedztwa wychodzącego.
\end{twierdzenie}
\begin{proof}
Jeżeli podczas przeglądania łuku $u\to v$
wierzchołek $v$ jest szary, należy on
do aktywnego stosu rekurencji i jest
przodkiem $u$. Łuk $u\to v$ wraz
z drzewową drogą skierowaną od $v$ do $u$
tworzy skierowany cykl.

Załóżmy odwrotnie, że DFS nie znalazło
żadnego łuku do szarego wierzchołka.
Rozważmy dowolny łuk $u\to v$.
W chwili jego przeglądania $v$ jest biały
albo czarny. Jeśli jest biały, wywołanie
\texttt{DFS(v)} skończy się przed zakończeniem
\texttt{DFS(u)}. Jeśli jest czarny, jego
wywołanie zakończyło się już wcześniej.
W obu przypadkach
$\mathrm{tout}[u]>\mathrm{tout}[v]$.
Istnienie skierowanego cyklu
$v_0\to v_1\to\cdots\to v_{k-1}\to v_0$
wymagałoby zatem nierówności
$\mathrm{tout}[v_0]>\mathrm{tout}[v_1]>
\cdots>\mathrm{tout}[v_{k-1}]>
\mathrm{tout}[v_0]$, co jest niemożliwe.
Otrzymujemy równoważność.

Dla acyklicznego $D$ stwierdzona
nierówność obowiązuje na każdym łuku,
zatem uporządkowanie wierzchołków
według malejących czasów wyjścia
jest porządkiem topologicznym.
Każdy wierzchołek i każdy łuk
analizujemy stałą liczbę razy,
więc koszt wynosi $O(|V|+|A|)$.
\end{proof}

\begin{przyklad}[Zależności pomiędzy czynnościami]
Niech wierzchołki $A,B,C,D$ oznaczają
czynności i niech łuki będą
$A\to C$, $B\to C$, $C\to D$.
Porządki $A,B,C,D$ oraz $B,A,C,D$
są topologiczne; $C,A,B,D$ nie jest,
ponieważ narusza zależność $A\to C$.
Wprowadzenie łuku $D\to A$
tworzy cykl $A\to C\to D\to A$,
więc porządek topologiczny przestaje istnieć.
\end{przyklad}

\section{Podsumowanie i dalsze algorytmy}
\label{sec:tg-algo-summary}

Dla reprezentacji listowej BFS i DFS mają jednakową
złożoność $O(n+m)$, ale dostarczają różnych danych:
BFS wyznacza najkrótsze drogi w grafie nieważonym
i warstwy metryczne; DFS dostarcza przedziały
czasowe, relacje przodków, mosty i porządki
topologiczne. W obu przypadkach złożoność
dotyczy \emph{całego} grafu także wtedy, gdy
jest niespójny, po uruchomieniu procedury
z kolejnych nieodkrytych korzeni.

W następnych uzupełnieniach algorytmicznych
rozpatrzymy najkrótsze drogi z wagami,
minimalne drzewa rozpinające oraz
maksymalny przepływ i minimalny przekrój.
Należy odróżniać wymienione zadania:
BFS minimalizuje liczbę krawędzi drogi,
a nie sumę dowolnych wag; zwykły DFS
nie wyznacza drzew o najmniejszym koszcie.
